\chapter{Metodologia}

Per a la realització d’aquest treball, hem decidit utilitzar una metodologia que ens permetés organitzar la feina de manera eficient i assegurar la qualitat del resultat final. La metodologia inclou tant l’organització del projecte i la manera de treballar amb el codi i les dades com la redacció de la memòria amb eines informàtiques adequades. La tria d’aquestes eines s’ha basat en la seva capacitat per permetre un treball ordenat, segur i proper a la manera en què treballen els professionals en projectes tecnològics i científics.

\section{Control de versions amb Git}

Una part fonamental de la metodologia ha estat l’ús de Git, que és un sistema de control de versions que ens ha permès guardar totes les modificacions que hem fet al projecte, registrar els canvis realitzats i poder recuperar versions anteriors quan ha estat necessari. Aquesta eina ens ha donat seguretat a l’hora de treballar amb el codi i amb els fitxers de suport, ja que ens ha permès provar noves funcionalitats en branques separades sense afectar la versió principal del projecte. Això ha estat especialment útil en el desenvolupament de l’algoritme per calcular les rutes i gestionar la informació de les estacions de Bicing i AMBici. D’aquesta manera hem pogut comprovar els resultats pas a pas i organitzar el projecte amb més eficàcia, mantenint un registre complet de l’evolució de la feina.

\section{Redacció del treball amb LaTeX}

Per a la redacció de la memòria hem utilitzat LaTeX, un sistema que permet elaborar documents amb un format uniforme i professional. Aquesta eina ha estat especialment útil per estructurar el treball en capítols i seccions, afegir figures i taules amb els peus numerats automàticament, generar índexos i bibliografia sense errors i obtenir un resultat final amb un aspecte net i coherent. Treballar amb LaTeX ens ha permès centrar-nos en el contingut i en l’explicació dels processos seguits sense preocupar-nos de l’estètica, ja que el sistema s’encarrega de donar coherència visual a tot el document.

\section{Organització i planificació del treball}

Un altre aspecte important de la metodologia ha estat planificar les diferents fases del treball de manera clara i ordenada. Primer hem realitzat una recerca prèvia sobre les dades disponibles de Bicing i AMBici, les xarxes de carrils bici i les eines per calcular rutes. A continuació hem establert un flux de treball que inclou la creació del prototip de l’aplicació, la implementació de l’algoritme per calcular rutes combinant ambdós serveis i la validació del seu funcionament amb casos reals. Aquesta planificació ha estat essencial per poder avançar de manera ordenada i evitar errors o repeticions.

\section{Avantatges de la metodologia aplicada}

La combinació de Git i LaTeX ha estat especialment beneficiosa perquè ens ha permès treballar de manera estructurada i segura. Hem pogut controlar l’evolució del projecte, tenir un registre de tots els canvis i redactar un document amb un format professional. Això ens ha acostat a la manera en què treballen els professionals en projectes tecnològics i científics, millorant la qualitat del treball i la comprensió dels processos seguits. També ha facilitat que el treball fos més transparent i fàcil de revisar, ja que tot està documentat i registrat de manera clara.

\section{Consideracions finals}

Gràcies a aquesta metodologia, hem aconseguit desenvolupar l’aplicació i redactar la memòria amb un alt grau d’eficiència i qualitat. El procés ha estat ordenat, segur i proper a la pràctica professional, cosa que aporta un valor afegit al treball de recerca. A més, aquesta manera de treballar ens ha permès aprendre a utilitzar eines que són habituals en l’àmbit científic i tecnològic, aportant-nos una experiència pràctica i formativa molt útil per a futurs projectes.
